\addcontentsline{toc}{chapter}{\textbf{Список использованных источников}}

\begin{thebibliography}{99}

\bibitem{tanenbaum2021}
Таненбаум Э., Ван Стеен М.  
Распределённые системы: принципы и парадигмы. — СПб.: Питер, 2021. — 928 с.

\bibitem{chernyshev2020}
Чернышев С. А., Кузнецов И. В.  
Распределённые системы и параллельные вычисления. — М.: Изд-во МГТУ им. Н. Э. Баумана, 2020. — 352 с.

\bibitem{sun2008ot}
Sun C., Ellis C.  
Operational Transformation in Real-Time Group Editors: Issues, Algorithms, and Achievements // *Computer Supported Cooperative Work (CSCW)*. — 2008. — Т. 17. — С. 69–100.

\bibitem{oster2006p2p}
Oster G., Urso P., Molli P., Imine A.  
Data Consistency for P2P Collaborative Editing // *Proceedings of the 2006 ACM Conference on Computer Supported Cooperative Work (CSCW’06)*. — 2006.

\bibitem{kleppmann2017crdt}
Kleppmann M., Beresford A. R.  
A Conflict-Free Replicated JSON Datatype // *IEEE Transactions on Parallel and Distributed Systems*. — 2017. — Vol. 28(10). — P. 2733–2746.

\bibitem{shevchenko2019protocols}
Шевченко А. А., Михайлов П. Н.  
Протоколы обмена сообщениями в распределённых системах. — М.: Лаборатория знаний, 2019. — 288 с.

\bibitem{vanderlinde2020automerge}
Van der Linde C., Kleppmann M.  
Automerge: A JSON-like Data Structure for Concurrent Editing. — (https://arxiv.org/abs/2004.00120), 2020.

\bibitem{google2019ot}
Google Inc.  
Operational Transformation Overview. — (https://developers.google.com/google-apps/realtime/overview) (дата обращения: 15.11.2025).

\bibitem{martin2020crdt}
Martin F., Oster G., Urso P.  
CRDTs: Consistency without Coordination // *Communications of the ACM*. — 2020. — Vol. 63(4). — P. 75–83.

\bibitem{rfc6455}
RFC 6455: The WebSocket Protocol. — IETF, 2011. — (https://datatracker.ietf.org/doc/html/rfc6455) (дата обращения: 12.11.2025).

\bibitem{yjs2025}
GitHub.  
Yjs — A CRDT Framework for Building Collaborative Applications. — (https://github.com/yjs/yjs) (дата обращения: 03.11.2025).

\bibitem{schmidt2022}
Шмидт В. Е.  
Проектирование и разработка клиент-серверных приложений. — М.: Академия, 2022. — 256 с.

\end{thebibliography}